\documentclass[12pt,a4paper]{article}

\usepackage[utf8]{inputenc}
\usepackage[T1]{fontenc}
\usepackage{amsmath,amssymb,amsthm}
\usepackage{booktabs}
\usepackage{graphicx}
\usepackage[margin=2.5cm]{geometry}
\usepackage{natbib}
\usepackage{hyperref}
\usepackage{cleveref}
\usepackage{setspace}
\usepackage{caption}
\usepackage{subcaption}
\usepackage{float}

\onehalfspacing

\title{External Shock Propagation to Hong Kong Under the Currency Board:\\
A VAR Analysis of US Monetary Policy and China Growth Spillovers}
\author{Steven Chung}
\date{\today}

\begin{document}
\maketitle

\begin{abstract}
How do US monetary policy shocks and China growth shocks propagate to
Hong Kong's real economy under the currency board? I estimate a
six-variable quarterly VAR for Hong Kong (1998Q1--2024Q3) with an
external-first Cholesky identification scheme. Forecast error variance
decomposition reveals a sharp division of influence: China GDP growth
explains 60--78\% of Hong Kong GDP forecast error variance at the
two-year horizon, while the US federal funds rate accounts for less than
5\%. The pattern reverses for interest rates, where the US FFR explains
59--75\% of HIBOR variance through the currency board peg. Historical
decomposition attributes the 2008--09 recession primarily to China
growth shocks and the post-2022 HIBOR spike to US monetary tightening.
Results are robust to alternative Cholesky orderings and Bayesian VAR
shrinkage.

\medskip\noindent
\textbf{Keywords:} Hong Kong, currency board, VAR, FEVD, monetary
transmission, China spillovers

\medskip\noindent
\textbf{JEL Classification:} E32, E52, F41, F42
\end{abstract}

\newpage
\tableofcontents
\newpage

% ==================================================================
\section{Introduction}
\label{sec:intro}
% ==================================================================

Hong Kong operates one of the world's most rigid exchange rate regimes:
a currency board pegging the Hong Kong dollar to the US dollar at a
narrow band since 1983. Under this arrangement, Hong Kong imports US
monetary policy directly---the Hong Kong Monetary Authority (HKMA)
adjusts the base rate in lockstep with the Federal Reserve. At the same
time, Hong Kong's real economy is deeply integrated with mainland China
through trade, tourism, and financial linkages. This creates a natural
tension: the interest rate channel is anchored to US conditions, while
the real activity channel is driven by Chinese demand.

This paper quantifies how US monetary policy shocks and China growth
shocks propagate to Hong Kong's domestic economy. Using a six-variable
quarterly VAR estimated on 1998Q1--2024Q3 data, I compute orthogonalized
impulse response functions, forecast error variance decompositions
(FEVD), and historical decompositions under a Cholesky identification
scheme that orders external variables first.

The main findings are:
\begin{enumerate}
    \item \textbf{China dominates real activity.} China GDP growth
    explains 60--78\% of Hong Kong GDP forecast error variance at the
    8-quarter horizon, depending on VAR vs.\ BVAR specification.
    The US federal funds rate explains less than 5\%.

    \item \textbf{The US dominates interest rates.} Through the currency
    board mechanism, the US FFR explains 59--75\% of 3-month HIBOR
    variance, while China growth shocks contribute less than 20\%.

    \item \textbf{Asymmetric transmission.} US monetary tightening
    raises HIBOR almost one-for-one but has modest direct effects on
    GDP and unemployment. China slowdowns, by contrast, depress GDP
    and raise unemployment without an equivalent interest rate offset.
\end{enumerate}

These results have direct policy relevance for HKMA's macroprudential
toolkit: since the interest rate channel is imported and the dominant
real shock source (China) cannot be offset by domestic monetary policy,
fiscal and prudential instruments must bear the stabilization burden.

The remainder of the paper is organized as follows.
\Cref{sec:literature} reviews related work on small open economy VARs
and Hong Kong macro modelling. \Cref{sec:data} describes the data and
variable construction. \Cref{sec:model} presents the econometric
framework. \Cref{sec:results} reports the main results.
\Cref{sec:robustness} checks robustness. \Cref{sec:conclusion}
concludes.

% ==================================================================
\section{Literature Review}
\label{sec:literature}
% ==================================================================

This paper contributes to three strands of literature.

\paragraph{Small open economy VARs.}
The structural VAR framework for analyzing external shock transmission
to small open economies originates with \citet{cushman1997identifying}
and \citet{kim2001international}. The standard approach places external
variables (world interest rates, commodity prices, foreign output) in a
block-exogenous position within the Cholesky ordering, an assumption
that is particularly natural for Hong Kong given the currency board's
automatic interest rate passthrough.

\paragraph{China--Hong Kong economic linkages.}
\citet{genberg2009monetary} study monetary transmission under the HKD
peg and find that US monetary conditions dominate short-term interest
rates but have limited direct effects on output. More recently,
\citet{he2014balancing} examine how mainland economic conditions
increasingly shape Hong Kong's business cycle. My contribution is to
provide a unified FEVD-based quantification of both channels
simultaneously.

\paragraph{Currency board macro dynamics.}
The macroeconomics of currency boards is surveyed by
\citet{hanke2002currency}. Under a strict currency board, the central
bank surrenders monetary autonomy in exchange for exchange rate
stability. The implication for VAR modelling is that the domestic
interest rate should be ordered after the anchor currency's policy
rate, reflecting its endogenous nature.

% ==================================================================
\section{Data}
\label{sec:data}
% ==================================================================

\subsection{Variables and Sources}

I use six quarterly variables spanning 1998Q1--2024Q3 (107 observations):

\begin{table}[H]
\centering
\caption{Variable Definitions and Sources}
\label{tab:variables}
\begin{tabular}{@{}llll@{}}
\toprule
Variable & Description & Unit & Source \\
\midrule
\texttt{us\_ffr} & US effective federal funds rate & \% & FRED \\
\texttt{china\_gdp} & China real GDP growth & YoY \% & World Bank \\
\texttt{gdp\_growth} & Hong Kong real GDP growth & YoY \% & World Bank \\
\texttt{cpi\_inflation} & Hong Kong CPI inflation & YoY \% & World Bank \\
\texttt{unemployment} & Hong Kong unemployment rate & \% & World Bank \\
\texttt{hibor\_3m} & 3-month HIBOR & \% & Derived (FFR + spread) \\
\bottomrule
\end{tabular}
\end{table}

The US federal funds rate is obtained as a monthly series from FRED
and aggregated to quarterly means. Hong Kong and China macro variables
are sourced from the World Bank's annual indicators and interpolated
to quarterly frequency using cubic spline temporal disaggregation.
HIBOR is proxied as the US FFR plus a calibrated spread, reflecting
the currency board's near-automatic passthrough.

\subsection{Stationarity and Cointegration}

I test each variable for stationarity using both ADF and KPSS tests.
Where both tests agree on a unit root, I take first differences. The
dual-test protocol resolves the known size-power trade-off between ADF
(conservative) and KPSS (liberal for trend-stationary processes).

Johansen cointegration testing on the I(1) subset
(\texttt{china\_gdp}, \texttt{cpi\_inflation}, \texttt{unemployment})
detects one cointegrating relationship at the 90\% level. I proceed
with the differenced VAR but note that a VECM specification may improve
long-horizon forecasts.

% ==================================================================
\section{Econometric Framework}
\label{sec:model}
% ==================================================================

\subsection{VAR Specification}

The reduced-form VAR(\(p\)) is:
\[
y_t = c + A_1 y_{t-1} + A_2 y_{t-2} + \cdots + A_p y_{t-p} + u_t,
\quad u_t \sim \mathcal{N}(0, \Sigma_u)
\]
where \(y_t\) is the \(K \times 1\) vector of endogenous variables
(\(K = 6\)), \(c\) is a vector of intercepts, \(A_\ell\) are
\(K \times K\) coefficient matrices, and \(u_t\) are reduced-form
residuals. Lag order \(p\) is selected by BIC subject to a
parameter-count guardrail (\(p = 3\)).

\subsection{Structural Identification}

I identify structural shocks via a Cholesky decomposition of
\(\Sigma_u = PP'\), where \(P\) is lower-triangular. The variable
ordering follows the external-first convention:
\[
[\text{us\_ffr}, \; \text{china\_gdp}, \; \text{gdp\_growth}, \;
\text{cpi\_inflation}, \; \text{unemployment}, \; \text{hibor\_3m}]
\]
This ordering assumes that (i) US and China variables do not respond
contemporaneously to Hong Kong shocks (small open economy assumption),
and (ii) domestic real variables (GDP, CPI) respond within the quarter
to interest rate shocks but not vice versa.

\subsection{Impulse Response Functions}

Orthogonalized IRFs are computed as:
\[
\Theta_h = \Phi_h P, \quad \text{where } \Phi_h = J F^h J',
\]
with \(F\) the companion matrix, \(J = [I_K \; 0 \; \cdots \; 0]\),
and \(P = \text{chol}(\Sigma_u)\).

\subsection{Forecast Error Variance Decomposition}

The FEVD at horizon \(h\) for variable \(i\) due to shock \(j\) is:
\[
\text{FEVD}_{i,j}(h) = \frac{\sum_{s=0}^{h} \theta_{ij,s}^2}
{\sum_{s=0}^{h} \sum_{k=1}^{K} \theta_{ik,s}^2}
\]
where \(\theta_{ij,s}\) is the \((i,j)\) element of \(\Theta_s\).
By construction, \(\sum_j \text{FEVD}_{i,j}(h) = 1\).

\subsection{Historical Decomposition}

I decompose each observed period into contributions from each
structural shock:
\[
y_t - \bar{y} = \sum_{j=1}^{K} \sum_{s=0}^{t-1} \Theta_s \,
e_{j,t-s}
\]
where \(e_t = P^{-1} u_t\) are the structural innovations and
\(\bar{y}\) is the deterministic baseline.

\subsection{Bayesian VAR}

As a robustness check, I estimate a Minnesota-prior BVAR. The prior
shrinks cross-variable coefficients toward zero and own-lag coefficients
toward a random walk:
\[
\text{Var}(A_{ij}^{(\ell)}) = \begin{cases}
(\lambda_1 / \ell^{\lambda_3})^2 & \text{if } i = j \\
(\lambda_1 \lambda_2 \hat\sigma_i / (\ell^{\lambda_3} \hat\sigma_j))^2
& \text{if } i \neq j
\end{cases}
\]
where \(\hat\sigma_i\) is the residual standard deviation from a
univariate AR for variable \(i\), and \(\lambda_1 = 0.2\),
\(\lambda_2 = 0.5\), \(\lambda_3 = 1.0\) are hyperparameters.

% ==================================================================
\section{Results}
\label{sec:results}
% ==================================================================

\subsection{FEVD: Who Drives What?}

\Cref{tab:fevd} presents the forecast error variance decomposition at
the 8-quarter horizon for both VAR and BVAR specifications.

\begin{table}[H]
\centering
\caption{Forecast Error Variance Decomposition at \(h = 8\) Quarters}
\label{tab:fevd}
\begin{tabular}{@{}lcccc@{}}
\toprule
& \multicolumn{2}{c}{VAR (BIC, $p=3$)} &
\multicolumn{2}{c}{BVAR (Minnesota, $p=3$)} \\
\cmidrule(lr){2-3} \cmidrule(lr){4-5}
Target & US FFR & China GDP & US FFR & China GDP \\
\midrule
HK GDP growth    & 3.0\% & 77.5\% & 4.1\% & 59.7\% \\
HK CPI inflation & 2.7\% & 28.7\% & 0.5\% & 7.8\% \\
HK Unemployment  & 9.2\% & 26.8\% & 5.8\% & 6.7\% \\
3M HIBOR         & 59.2\% & 18.9\% & 75.2\% & 0.9\% \\
\bottomrule
\end{tabular}
\end{table}

The central finding is a sharp division of influence:
\begin{itemize}
    \item China GDP growth is the dominant driver of Hong Kong's real
    activity (60--78\% of GDP variance). The US federal funds rate
    contributes less than 5\%.

    \item The pattern reverses for interest rates: the US FFR explains
    59--75\% of HIBOR variance, confirming the currency board's
    mechanical passthrough.

    \item CPI and unemployment receive mixed influence from both
    sources, with China growth mattering more in the OLS VAR
    specification.
\end{itemize}

\subsection{Historical Decomposition}

The historical decomposition identifies which shocks drove specific
episodes. During the 2008--09 Global Financial Crisis, the China growth
shock (negative) accounts for the largest share of Hong Kong's GDP
decline. During 2022--23, the US monetary tightening shock dominates
the rise in HIBOR, while China's post-COVID slowdown drives the real
economy weakness.

\subsection{Forecast Evaluation}

Expanding-window out-of-sample evaluation (minimum training window: 60
quarters, horizons: 1Q and 4Q) compares the VAR against AR(1) and
random walk benchmarks. At the 1-quarter horizon, the VAR outperforms
benchmarks for GDP growth and unemployment. At 4 quarters, no model
consistently dominates, consistent with the difficulty of long-horizon
macro forecasting.

% ==================================================================
\section{Robustness}
\label{sec:robustness}
% ==================================================================

\subsection{Cholesky Ordering}

I re-estimate the model under three orderings: (i) the baseline
external-first ordering, (ii) domestic variables first, and (iii)
reversed. The qualitative result---China dominates real activity,
US dominates HIBOR---is preserved across all orderings, though the
exact FEVD shares shift by 5--15 percentage points.

\subsection{Sub-sample Stability}

I estimate on pre-GFC (1998--2009), post-GFC (2010--2024), and
full-sample-excluding-COVID sub-samples. The full sample and
ex-COVID models are stable (max eigenvalue $< 1$). The post-GFC
sub-sample shows mild instability, likely reflecting the shorter
sample combined with the 2022--23 rate hiking cycle.

\subsection{VAR vs.\ BVAR}

The BVAR (Minnesota prior, \(\lambda_1 = 0.2\)) yields qualitatively
similar FEVD patterns with tighter shrinkage. The BVAR is always
stable (max eigenvalue 0.91 vs.\ 0.99 for OLS VAR) and shows lower
coefficient norms, suggesting less sensitivity to outlier observations.

% ==================================================================
\section{Conclusion}
\label{sec:conclusion}
% ==================================================================

This paper provides a quantitative answer to how US and China shocks
propagate to Hong Kong under the currency board. The evidence shows a
clean separation: China growth dominates the real economy, while US
monetary policy dominates the interest rate channel. This asymmetry
has direct implications for macroprudential policy design---since the
currency board precludes independent monetary responses to real shocks,
fiscal and regulatory instruments must bear the stabilization burden
when China's economy slows.

Several avenues for future work remain. First, real quarterly data from
official Hong Kong sources (C\&SD, HKMA) would strengthen the empirical
foundation. Second, a VECM specification could exploit the detected
cointegrating relationship for improved long-horizon forecasts. Third,
time-varying parameter VARs could capture the evolving relative
importance of US and China shocks as mainland financial integration
deepens.

% ==================================================================
\bibliographystyle{apalike}
\bibliography{references}
% ==================================================================

\end{document}
